\begin{objectivebox}[이 장에서 붙잡을 질문]
가능한 입력을 빠짐없이 적으면, 우리가 원하는 규칙은 어떻게 하나의 함수가 될까?
\end{objectivebox}

\section{두 조건을 빠짐없이 적는 법}

작업실에 작은 절단기가 있다고 해 보자. 안전 덮개가 닫혀 있고 작업자가 시작 단추를 눌렀을 때만 모터가 돌아야 한다. 두 조건 가운데 하나라도 만족하지 않으면 모터는 멈춰 있어야 한다.

이것은 논리 규칙을 설명하기 위한 단순한 모형이다. 실제 기계의 안전장치는 센서 고장과 단선, 비상 정지처럼 더 많은 상황을 고려해 훨씬 조심스럽게 설계한다. 여기서는 두 조건만 남겨, 문장으로 적은 규칙을 계산 가능한 형태로 바꾸는 과정에 집중한다.

먼저 조건이 맞으면 \(1\), 맞지 않으면 \(0\)이라고 쓰자.

\begin{description}[style=nextline,leftmargin=2.5em,font=\sffamily\bfseries]
  \item[입력 \(A\)] 안전 덮개가 닫혔는가?
  \item[입력 \(B\)] 시작 단추를 눌렀는가?
  \item[출력 \(Y\)] 모터를 켜도 되는가?
\end{description}

문장만 읽으면 빠뜨린 경우가 없는지 한눈에 확인하기 어렵다. 덮개는 닫혔지만 단추를 누르지 않은 경우는 어떻게 되는가? 덮개가 열린 채로 단추를 누르면 어떻게 되는가? 두 질문을 포함해 가능한 조합을 모두 적으면 다음 표가 된다.

\begin{center}
\small
\begin{tabular}{ccc}
\toprule
\sffamily\bfseries \(A\): 덮개 & \sffamily\bfseries \(B\): 단추 & \sffamily\bfseries \(Y\): 모터 허가 \\
\midrule
0 & 0 & 0 \\
0 & 1 & 0 \\
1 & 0 & 0 \\
1 & 1 & 1 \\
\bottomrule
\end{tabular}
\end{center}

첫째 행에서는 덮개도 열려 있고 단추도 누르지 않았으므로 모터를 켜지 않는다. 둘째 행과 셋째 행에서는 조건 하나만 맞으므로 역시 모터를 켜지 않는다. 마지막 행에서만 두 조건이 모두 맞으므로 출력이 \(1\)이다.

\begin{bookfigure}{일상 규칙의 모든 경우를 적으면 진리표가 된다}{fig:ch01-rule-to-truth-table}
\begin{tikzpicture}[diagram]
  \path[use as bounding box] (0,0) rectangle (13.0,7.75);

  \draw[diagram panel,fill=black!7] (0.05,0.42) rectangle (4.18,7.62);
  \node[diagram panel title] at (0.30,7.27) {단순화한 규칙};
  \node[diagram note,anchor=west] at (0.35,6.55) {안전 덮개가 닫혀 있고};
  \node[diagram note,anchor=west] at (0.35,6.02) {시작 단추가 눌렸을 때만};
  \node[diagram note,anchor=west] at (0.35,5.49) {모터 작동을 허가한다.};

  \node[concept box,minimum width=31mm,minimum height=16mm] at (2.12,3.82)
    {둘 다 1일 때만\\\(Y=1\)};
  \node[diagram note,anchor=west,text=black!68] at (0.35,2.65) {\(A\) = 안전 덮개};
  \node[diagram note,anchor=west,text=black!68] at (0.35,2.13) {\(B\) = 시작 단추};
  \node[diagram note,anchor=west,text=black!68] at (0.35,1.61) {\(Y\) = 모터 허가};
  \node[diagram note,align=center,text=black!62] at (2.12,0.86)
    {실제 안전회로가 아닌\\교육용 모형};

  \draw[concept arrow] (4.35,4.05) -- (5.02,4.05);

  \draw[diagram panel] (5.20,0.42) rectangle (12.95,7.62);
  \node[diagram panel title] at (5.47,7.27) {질문 순서대로 네 경우를 채운다};

  \node[concept box,minimum width=14mm,minimum height=8mm] at (6.32,5.98) {\(A\)};
  \node[concept box,minimum width=14mm,minimum height=8mm] at (8.18,5.98) {\(B\)};
  \node[concept emphasis,minimum width=29mm,minimum height=8mm] at (10.92,5.98) {\(Y=A\ \mathrm{AND}\ B\)};

  \foreach \y/\a/\b/\out/\num in {
    5.04/0/0/0/1,
    4.18/0/1/0/2,
    3.32/1/0/0/3,
    2.46/1/1/1/4
  }{
    \node[diagram note,anchor=east,text=black!58] at (5.72,\y) {\num};
    \node[gate,minimum width=14mm] at (6.32,\y) {\a};
    \node[gate,minimum width=14mm] at (8.18,\y) {\b};
    \node[path gate,minimum width=29mm] at (10.92,\y) {\out};
  }

  \node[concept emphasis,minimum width=64mm,minimum height=14mm] at (9.08,1.15)
    {입력 2개 \(\longrightarrow 2^2=4\)행\\
     네 행을 모두 적으면 입력 조합을 빠뜨리지 않는다};
\end{tikzpicture}
\end{bookfigure}


\figref{fig:ch01-rule-to-truth-table}처럼 가능한 입력 조합을 모두 나열하고 각 경우의 출력을 적은 표를 \keyterm{진리표}(truth table)라고 한다. 이름은 거창하지만 하는 일은 규칙표와 같다. 중요한 점은 어떤 경우도 짐작으로 건너뛰지 않는다는 것이다. 진리표는 유한한 개수의 \(0\)과 \(1\)을 입력받는 논리 기능을 완전하게 적는 한 방법이다\parencite{mit6004truth}.

\section{빠짐없는 표는 얼마나 커지는가}

입력 \(A\) 하나만 있다면 가능한 값은 \(0\)과 \(1\), 두 가지다. 입력 \(B\)가 하나 더 생기면 \(A\)의 각 경우마다 \(B\)의 두 경우를 붙여야 한다. 따라서 가능한 조합의 수는

\[
  2\times2=2^2=4
\]

가 된다. 입력이 세 개라면 각 행이 다시 두 갈래로 나뉘므로

\[
  2\times2\times2=2^3=8
\]

행이 필요하다. 일반적으로 입력이 \(n\)개이고 각 입력이 \(0\) 또는 \(1\)이라면 가능한 입력 조합은 \(2^n\)개다.

이 숫자는 진리표의 장점과 대가를 동시에 보여 준다. 모든 행을 쓰면 빠진 경우가 없다는 사실을 직접 확인할 수 있다. 그러나 입력이 하나 늘 때마다 행의 수가 두 배가 된다. 입력 열 개에는 \(2^{10}=1{,}024\)행, 스무 개에는 \(2^{20}=1{,}048{,}576\)행이 필요하다. 가능한 경우를 모두 적는 방법은 정확하지만, 입력이 많아지면 표 자체가 빠르게 커진다.

여기서 처음으로 계산이 요구하는 양을 셀 수 있다. 두 입력 규칙을 표로 완전히 적으려면 네 행이 필요하다. 아직 회로가 몇 개인지, 시간이 얼마나 걸리는지는 세지 않았다. 네 행은 규칙의 크기이지 기계의 크기가 아니다.

\section{진리표는 부품이 아니다}

진리표 자체가 모터를 켜는 것은 아니다. 표에는 전선도 없고 전류도 흐르지 않는다. 진리표는 회로가 지켜야 할 규칙을 적은 \keyterm{명세}(specification)다. 표는 입력마다 어떤 출력이 나와야 하는지를 정하지만, 그것을 직접 만들지는 않는다.

방금 적은 규칙은 두 입력이 모두 \(1\)일 때만 출력이 \(1\)이다. 이 입력과 출력의 관계를 \keyterm{AND}(논리곱) 함수라고 부른다. 함수라는 말은 여기서 거대한 수식을 뜻하지 않는다. 입력을 하나 정하면 그 입력에 대응하는 출력이 정해진다는 뜻이다. 네 행 전체가 AND 함수의 약속이다.

같은 약속은 릴레이나 트랜지스터 논리 게이트처럼 여러 물리적 방법으로 실현할 수 있다. 이 책에서는 그 세부를 잠시 접고, \(0\)과 \(1\)을 입력받아 \(0\)과 \(1\)을 출력하는 디지털 논리 수준에서 시작한다. 무엇을 해야 하는지 정하는 명세와 그것을 어떻게 만드는지 정하는 구현을 구분해야, 같은 답을 내는 여러 회로의 비용을 비교할 수 있다.

\section{NAND는 AND의 결과를 뒤집는다}

진리표마다 전혀 다른 종류의 부품을 새로 만들어야 한다면 필요한 부품의 종류가 끝없이 늘어난다. 반대로 한 종류의 단순한 부품을 여러 방식으로 연결해 다양한 규칙을 만들 수 있다면, 복잡한 계산을 같은 재료의 조합으로 설명할 수 있다. 그런 출발점으로 사용할 수 있는 논리 게이트 가운데 하나가 NAND다\parencite{nand2tetrisproject1}.

\keyterm{NAND}(Not AND, 부정 논리곱)는 AND의 결과를 뒤집는다. 두 입력이 모두 \(1\)일 때 AND의 출력은 \(1\)이므로 NAND의 출력은 \(0\)이다. 나머지 세 경우에는 AND가 \(0\)이므로 NAND는 \(1\)을 낸다.

\begin{center}
\small
\begin{tabular}{cccc}
\toprule
\sffamily\bfseries \(A\) & \sffamily\bfseries \(B\) & \sffamily\bfseries AND(\(A,B\)) & \sffamily\bfseries NAND(\(A,B\)) \\
\midrule
0 & 0 & 0 & 1 \\
0 & 1 & 0 & 1 \\
1 & 0 & 0 & 1 \\
1 & 1 & 1 & 0 \\
\bottomrule
\end{tabular}
\end{center}

표와 부품은 여기서도 구분해야 한다. NAND 진리표는 게이트가 지켜야 할 입력과 출력의 관계다. NAND 게이트는 그 관계를 실제로 구현하는 논리 부품이다. 이 책에서 게이트 상자 하나를 그릴 때는 내부 트랜지스터를 세는 대신 그 상자가 이 진리표를 따른다고 약속한다.

그런데 우리가 원한 모터 허가 규칙은 NAND가 아니라 AND였다. 한 종류의 부품만 사용한다는 선택 아래에서 NAND로 AND와 똑같은 마지막 열을 만들 수 있을까?

첫 번째 NAND 게이트에 \(A\)와 \(B\)를 넣고 그 출력을 \(N\)이라고 하자.

\[
  N=\operatorname{NAND}(A,B)
\]

두 번째 NAND 게이트의 두 입력에는 똑같이 \(N\)을 넣는다.

\[
  Y=\operatorname{NAND}(N,N)
\]

입력이 \(N=0\)이면 NAND의 두 입력은 \(0,0\)이므로 출력은 \(1\)이다. 입력이 \(N=1\)이면 두 입력은 \(1,1\)이므로 출력은 \(0\)이다. 같은 신호를 두 입력에 넣은 NAND는 그 값을 반대로 바꾸는 장치처럼 작동한다. 입력 하나의 값을 반대로 만드는 이 기능을 \keyterm{NOT}(부정)이라고 한다. 따라서 두 번째 게이트는 첫 번째 게이트의 출력을 한 번 더 뒤집는다.

중간값까지 표로 확인하면 다음과 같다.

\begin{center}
\small
\begin{tabular}{cccc}
\toprule
\sffamily\bfseries \(A\) & \sffamily\bfseries \(B\) & \sffamily\bfseries \(N=\operatorname{NAND}(A,B)\) & \sffamily\bfseries \(Y=\operatorname{NAND}(N,N)\) \\
\midrule
0 & 0 & 1 & 0 \\
0 & 1 & 1 & 0 \\
1 & 0 & 1 & 0 \\
1 & 1 & 0 & 1 \\
\bottomrule
\end{tabular}
\end{center}

마지막 열은 처음 만든 모터 허가 진리표와 정확히 같다. NAND 게이트 두 개를 연결해 AND 함수를 구현한 것이다.

\begin{bookfigure}{진리표는 명세이고, 게이트 연결은 구현이다}{fig:ch01-specification-vs-implementation}
\begin{tikzpicture}[diagram]
  \path[use as bounding box] (0,0) rectangle (13.0,7.85);

  \draw[diagram panel] (0.05,0.75) rectangle (4.10,7.70);
  \node[diagram panel title] at (0.30,7.34) {명세 · 무엇을 해야 하는가};

  \node[concept box,minimum width=10mm] at (0.92,6.35) {\(A\)};
  \node[concept box,minimum width=10mm] at (2.05,6.35) {\(B\)};
  \node[concept emphasis,minimum width=10mm] at (3.20,6.35) {\(Y\)};
  \foreach \y/\a/\b/\out in {
    5.43/0/0/0,
    4.57/0/1/0,
    3.71/1/0/0,
    2.85/1/1/1
  }{
    \node[gate,minimum width=10mm] at (0.92,\y) {\a};
    \node[gate,minimum width=10mm] at (2.05,\y) {\b};
    \node[path gate,minimum width=10mm] at (3.20,\y) {\out};
  }
  \node[diagram note,align=center] at (2.05,1.77) {AND의 네 입력·출력 관계};
  \node[diagram note,align=center,text=black!62] at (2.05,1.20) {지켜야 할 마지막 열};

  \draw[concept dashed arrow] (4.27,4.42) -- (5.10,4.42);
  \node[diagram note,align=center] at (4.68,3.77) {같은\\마지막 열};

  \draw[diagram panel] (5.28,0.75) rectangle (12.95,7.70);
  \node[diagram panel title] at (5.55,7.34) {구현 · 어떻게 만들 것인가};

  \node[diagram note,anchor=east] (ain) at (5.96,5.54) {\(A\)};
  \node[diagram note,anchor=east] (bin) at (5.96,3.54) {\(B\)};
  \node[path gate,minimum width=22mm,minimum height=14mm] (nand1) at (7.45,4.54)
    {NAND 1\\\(N=\operatorname{NAND}(A,B)\)};
  \node[path gate,minimum width=23mm,minimum height=14mm] (nand2) at (10.65,4.54)
    {NAND 2\\\(\operatorname{NOT}(N)\)};
  \node[diagram note,anchor=west] (yout) at (12.43,4.54) {\(Y\)};

  \draw[path wire] (ain.east) -| (nand1.north west);
  \draw[path wire] (bin.east) -| (nand1.south west);
  \coordinate (branch) at (8.98,4.54);
  \draw[path wire] (nand1.east) -- node[above,font=\DiagramBodyFont] {\(N\)} (branch);
  \fill (branch) circle (1.15pt);
  \draw[path wire] (branch) |- ($(nand2.west)+(0,0.25)$);
  \draw[path wire] (branch) |- ($(nand2.west)+(0,-0.25)$);
  \draw[path wire] (nand2.east) -- (yout.west);

  \node[concept emphasis,minimum width=56mm,minimum height=10mm] at (9.15,2.03)
    {NAND 2개 · 가장 긴 경로 2층};
  \node[diagram note,align=center,text=black!65] at (9.15,1.22)
    {한 종류의 게이트로 AND의 마지막 열을 되찾는다};

  \draw[gate group] (0.10,0.12) rectangle (12.92,0.60);
  \node[diagram footer] at (6.50,0.36)
    {접어 둔 층: 전압 · 트랜지스터 · 공정};
\end{tikzpicture}
\end{bookfigure}


\figref{fig:ch01-specification-vs-implementation}의 왼쪽은 보존해야 할 네 행의 명세이고, 오른쪽은 그 명세를 만족시키는 한 가지 구현이다. NAND 게이트는 두 개다. 입력 \(A\)나 \(B\)의 변화가 출력 \(Y\)에 이르기까지 지나가는 가장 긴 경로도 NAND 두 층이다. 이 숫자는 지금 그린 교육용 회로의 상자와 경로를 센 값이다. 실제 칩의 면적이나 실행 시간을 직접 잰 값은 아니다.

한 종류의 게이트로 조합할 수 있다는 단순성을 얻는 대신, AND를 하나의 기본 게이트로 허용할 때보다 게이트 수와 신호 경로가 늘어났다. 이것이 이 책에서 처음 만나는 설계의 교환이다. 같은 마지막 열을 지켰지만 구현의 비용은 달라질 수 있다.

\section{0과 1에서는 논리곱과 숫자 곱이 닮았다}

AND의 네 출력을 \(0\)과 \(1\)의 곱셈표와 나란히 놓아 보자.

\begin{center}
\small
\begin{tabular}{cccc}
\toprule
\sffamily\bfseries \(A\) & \sffamily\bfseries \(B\) & \sffamily\bfseries AND(\(A,B\)) & \sffamily\bfseries \(A\times B\) \\
\midrule
0 & 0 & 0 & 0 \\
0 & 1 & 0 & 0 \\
1 & 0 & 0 & 0 \\
1 & 1 & 1 & 1 \\
\bottomrule
\end{tabular}
\end{center}

\(0\)과 \(1\) 한 자리에서는 AND와 곱셈의 출력이 정확히 같다. 그렇다고 AND 게이트 하나가 모든 숫자의 곱셈을 곧바로 해 준다는 뜻은 아니다. 여러 자리 이진수에는 각 자리의 위치가 나타내는 값이 있고, 부분 결과를 더하는 과정도 필요하다. 우리는 3장에서 \(0\)과 \(1\)에 자릿값을 부여한 뒤, 5장에서 이 작은 표가 부분곱을 만드는 데 어떻게 쓰이는지 확인할 것이다.

지금까지 오른 첫 계단은 세 층으로 정리할 수 있다.

\begin{enumerate}
  \item 진리표로 입력과 출력의 규칙을 명세한다.
  \item 게이트로 그 규칙을 구현한다.
  \item 게이트를 조합해 처음에 기본 부품으로 주어지지 않은 기능을 만든다.
\end{enumerate}

실제 칩의 \(0\)과 \(1\)은 종이에 적힌 숫자가 아니라 정해진 범위의 전압으로 표현된다. NAND 게이트도 트랜지스터로 구현된다. 입력이 바뀐 뒤 출력이 바뀌기까지는 아주 짧지만 \(0\)이 아닌 시간이 걸리며, 게이트와 배선은 면적과 에너지를 사용한다. 이 장에서는 그 아래층을 접어 두었다. 같은 입력과 출력의 관계를 보존한다면 내부의 물리적 세부를 잠시 감추고 더 큰 기능을 조립할 수 있다는 사실이 지금은 더 중요하다.

다음 장에서는 같은 XOR 진리표를 만드는 두 NAND 회로를 나란히 놓는다. 그때부터 질문은 ``만들 수 있는가''에서 ``같은 함수를 어떤 식과 회로로 만드는 편이 더 적고 얕은가''로 바뀐다.

\section*{이번 장의 선택과 비용}
\addcontentsline{toc}{section}{이번 장의 선택과 비용}

\begin{center}
\small
\begin{tabular}{>{\raggedright\arraybackslash\sffamily\bfseries}p{0.29\textwidth} >{\raggedright\arraybackslash}p{0.58\textwidth}}
\toprule
보존한 기능·정확한 결과 & AND의 네 입력·출력 관계 \(0,0,0,1\) \\
\midrule
계산이 요구하는 양 & 빠짐없는 입력 조합 4행, AND를 NAND로 만들 때 NAND 기능 적용 2회 \\
\midrule
기계가 제공할 자원 & NAND 게이트 2개와 연결 통로, 입력에서 출력까지 가장 긴 NAND 경로 2층 \\
\midrule
설계 대가 & 한 종류의 게이트로 조합하는 단순성을 얻는 대신 AND를 기본 게이트로 허용할 때보다 게이트 수와 경로가 늘어남 \\
\midrule
아직 해결되지 않은 제약 & 같은 함수를 만드는 여러 식과 NAND 회로 가운데 무엇이 더 적고 얕은지 비교하는 기준 \\
\bottomrule
\end{tabular}
\end{center}
